\begin{abstract}
Quantum error mitigation (QEM) is vital for achieving utility-scale quantum computing on noisy intermediate-scale quantum devices. While learning-based approaches offer flexible mitigation capabilities, traditional architectures often struggle to efficiently generalize across varying circuit depths and fail to explicitly decouple distinct noise processes. In this work, we present a novel, physics-informed generative framework for quantum error mitigation that specifically targets coherent noise accumulation. Our method introduces a probabilistic generative process that models local coherent errors as latent space representations of noise generators, which are subsequently propagated through a deterministic module derived from physical principles. This architecture extends the original Neural Noise Accumulation Surrogate (NNAS) framework proposed in \cite{xu2025} by incorporating structurally distinct surrogate branches for independent stochastic and coherent error tracking, alongside a regularized loss function that latent's distributions with the prior conditional one. To validate this framework, we developed the entire software pipeline from the ground up, encompassing the models, the training routines, and the mixed-noise synthetic dataset generation. Numerical evaluations provide proof-of-concept results showing that integrating a structured, physics-informed coherent branch successfully captures complex error characteristics without degrading the underlying stochastic mitigation performance of the surrogate framework, offering a principled foundation for more thorough research.
\end{abstract}